% !PRINTED 51
\weberpage{51}

It is perhaps desirable to illustrate these theorems by a simple
example. Let the degree of the group $S$ be the square of a
prime $n = p^2$. Then $S$ has either only the one
invariant $p^2$, and is then cyclic:
\[
\text{a)}\quad S = 1, A, A^2 \dots A^{p^2-1},
\]
or $S$ has two invariants, which are both equal to $p$; then
$S$ consists of the elements:
\[
\text{b)}\quad S = A_1^{\alpha_1}A_2^{\alpha_2}, \quad \alpha_1,
\alpha_2 = 0, 1 \dots p-1.
\]
In case a) we obtain the $p^2$ characters
\[
\chi_B(A^\alpha) = \varepsilon^{\beta\alpha},
\]
if $\varepsilon$ is a primitive root of the equation
$\varepsilon^{pp} = 1$.
% [sic: the source prints \varepsilon^{pp}; the cyclic group order and surrounding p^2 formulas require \varepsilon^{p^2}.]
If, in order to form the divisors of $S$ according to Theorem 10.,
we take one of the characters $\chi_B$ in which
$\beta$ is not divisible by $p$, then $\chi_B(A^\alpha)$
can be $= 1$ only when $\alpha$ is divisible by $p^2$,
that is, $= 0$; i.\,e., we obtain only the divisor of $S$ consisting
of the unit element. But if we assume
$\beta = p\beta'$ to be divisible by $p$, but not by $p^2$, then
$\chi_{p\beta'}(A^\alpha)$ is $= 1$ then and only then when
$\alpha = p\alpha'$ is divisible by $p$. We thus obtain the
divisor
\[
T = 1, A^p, A^{2p} \dots A^{(p-1)p},
\]
and the divisor $U$ reciprocal to $T$ is here identical with $T$.

In case b) we must, in order to form the characters, assume two
$p$th roots of unity $\varepsilon^{\beta_1},
\varepsilon^{\beta_2}$, and obtain, if
\[
B = A_1^{\beta_1}A_2^{\beta_2}
\]
is set,
\[
\chi_B(A_1^{\alpha_1}A_2^{\alpha_2}) = \varepsilon^{\alpha_1\beta_1 + \alpha_2\beta_2}.
\]
If we set two of these characters $= 1$, that is
\[
\alpha_1\beta_1 + \alpha_2\beta_2 \equiv 0, \quad
\alpha_1\beta_1' + \alpha_2\beta_2' \equiv 0,
\]
then it follows, if the determinant $\beta_1\beta_2' - \beta_2\beta_1'$
is not divisible by $p$, that $\alpha_1$ and $\alpha_2$ must be
divisible by $p$. But if the determinant is divisible by $p$,
then the one of these two congruences is a consequence of the
other. In the first case, therefore, we obtain a group consisting
only of the unit element. We therefore obtain all divisors
$T_{\beta_1,\beta_2}$ distinct from 1 and $S$ when, for a fixed
$\beta_1, \beta_2$, we make the $\alpha_1, \alpha_2$ satisfy in all
possible ways the congruence
\weq{9}{\alpha_1\beta_1 + \alpha_2\beta_2 \equiv 0\ (\mathrm{mod}\ p)}
%CONT>>
